Feromagnetismus vyžaduje \textbf{nevykompenzované magnetické momenty atomů}, zejména spiny elektronů v neúplně obsazených vnitřních hladinách, a \textbf{výměnné síly} mezi sousedními atomy. Ty zajišťují spontánní uspořádání momentů v doménách.
\begin{itemize}
\item \textbf{Fero-:} momenty jsou souhlasně orientované, výsledná spontánní magnetizace je nenulová; Fe, Co, Ni.
\item \textbf{Antifero-:} podmřížky mají opačně orientované, stejně velké magnetizace, které se vyruší; např. MnO.
\item \textbf{Feri-:} podmřížky mají opačně orientované, ale nestejně velké magnetizace, takže se nevyruší; např. magnetit Fe$_3$O$_4$.
\end{itemize}
